\chapter{Exercise solutions and mentor notes}

These are reference solutions, not the only valid wording. The important quality
is an explicit boundary, justified protocol rule, and observable evidence. A
mentor should ask the intern to explain differences before showing this appendix.

\section{Foundations}

\subsection*{Exercise 1.1 --- Draw the journey}

For \url{http://127.0.0.1:8080/health}, no DNS lookup is required because the
host is already an address. IP routing selects loopback; packets do not leave the
machine. There is a TCP handshake, no TLS handshake, then HTTP request parsing,
routing to health, response serialization, and TCP close or reuse. Through an
HTTPS reverse proxy, add DNS, a route across one or more networks, TLS termination
at the proxy, one HTTP connection from client to proxy, and usually another
proxy-to-application connection. Mark the proxy as a trust boundary.

\subsection*{Exercise 1.2 --- Classify failures}

``Name not known'' belongs to DNS. ``Connection refused'' is TCP evidence that
the destination actively rejected the selected endpoint, commonly because no
listener exists. Certificate mismatch is TLS identity verification. A valid 400
is HTTP/application parsing evidence. ``Database busy'' occurs after enough HTTP
and routing succeeded to reach storage. Confirm with resolver output, the exact
endpoint, TLS verification detail, raw response, and structured server logs.

\subsection*{Exercise 2.1 --- Design a frame}

One valid design is a four-byte unsigned big-endian length followed by that many
UTF-8 bytes, repeated until the connection closes. Length zero represents an
empty note; close between frames ends the sequence; close before the declared
length is an incomplete-note error. Reject lengths above a configured maximum
before allocating. A delimiter design is possible, but then delimiter escaping
or prohibition must be specified.

\subsection*{Exercise 2.2 --- Count the bytes}

\code{hello} is five characters and five UTF-8 bytes. \code{café} is four
characters and five UTF-8 bytes when the last character is the single code point
U+00E9. A visually identical decomposed form can have five characters/code points
and six bytes. Network length is the length of the final encoded bytes, not the
visual character count.

\subsection*{Exercise 3.1 --- Endpoint inventory}

Example: listener \code{127.0.0.1:8080}; first client
\code{127.0.0.1:53001}; TCP tuple client-to-server is
\code{127.0.0.1:53001 -> 127.0.0.1:8080}. A second client might use port 53002.
Protocol and server endpoint remain; the client source port changes, allowing the
operating system to distinguish connections.

\subsection*{Exercise 3.2 --- Resolution narrative}

The stub asks a recursive resolver. On a miss, that resolver asks root for a
delegation, a TLD server for the authoritative delegation, and the authoritative
server for the record; it caches referrals and answer under TTLs. On a complete
cache hit, the resolver answers immediately and all upstream DNS round trips
disappear. The client may also have a local cache, removing even the resolver
exchange.

\subsection*{Exercise 4.1 --- Stream outcomes}

Legal reads include 800 once; 600 then 200; 100 repeated eight times; or 50,
300, 7, 443. TCP can split anywhere. The CRLF--CRLF delimiter ends fields, then
Content-Length (200 in this example) determines body completion. The header's
physical segment boundaries are irrelevant.

\subsection*{Exercise 4.2 --- Slow client budget}

Sixteen clients can each occupy one worker while never completing a header; the
seventeenth remains queued or unaccepted. A reasonable learning policy is a
five-second read timeout, 64-slot queue, and best-effort 503 on overflow. Shorter
timeouts free capacity faster but reject slower legitimate clients; larger queues
absorb bursts but consume descriptors and increase waiting latency.

\subsection*{Exercise 5.1 --- Method audit}

POST creates; GET reads; PUT replaces; PATCH partially modifies; DELETE removes.
GET, PUT, and DELETE are idempotent by standard semantics and are candidates for
retry after uncertain transport failure, though preconditions and concurrent
changes matter. POST creation is not safely retried without an idempotency key.
PATCH safety depends on the patch semantics: assignment of a fixed value is
idempotent, while ``increment by one'' is not.

\subsection*{Exercise 5.2 --- Choose the result}

Creation: 201 with Location. Successful no-body deletion: 204. Malformed JSON:
400. Unknown task: 404. Known path/wrong method: 405 plus Allow. Unsupported
media: 415. Invalid fields: 422 with field errors. Unexpected database failure:
500 with generic client detail and internal logging.

\subsection*{Exercise 6.1 --- Frame two requests}

A GET ends at the second CRLF of its empty line; the next method byte follows
immediately. A POST boundary is \code{header\_end + Content-Length}. There is no
extra blank line after the body. A parser must retain the next method bytes in
its connection buffer.

\subsection*{Exercise 6.2 --- Ambiguity review}

If a proxy honors Transfer-Encoding while an origin honors Content-Length, the
proxy can treat chunk syntax as the first body while the origin consumes four
raw bytes. Their next-request positions differ. Attacker-controlled trailing
bytes may become a second request for one recipient. The correct project policy
is reject the combination and close.

\subsection*{Exercise 7.1 --- Theory gate}

A complete diagram includes URL parsing, DNS, routed IP, TCP, optional TLS at a
proxy, HTTP bytes, parser limits/framing, request object, router parameters,
JSON/domain validation, parameterized SQL transaction, response object, HTTP
serialization, and the reverse byte path. Mark client/proxy and HTTP/application
boundaries as untrusted; mark bytes-to-text conversions at request line/fields,
percent-decoded target, JSON UTF-8, and response JSON.

\subsection*{Exercise 7.2 --- Feature classification}

Implemented: Content-Length, GET, ordinary repeated fields, JSON, and bounded
keep-alive. Rejected: HTTP/2 at the HTTP/1 parser, chunked/any Transfer-Encoding,
duplicate Content-Length, and unsupported methods/versions. Explanation-only:
TLS and authentication. Rejection is a response/closed connection or a separate
deployment boundary, never silent guessing.

\section{Building the server}

\subsection*{Exercise 8.1 --- Predict before running}

The stage answers every received byte sequence with the same 200 because it does
not parse. A changed path and POST therefore still succeed. Sending a line is
enough to unblock recv. Connecting and sending nothing blocks forever. A second
client may complete its handshake and wait in the listen backlog, but the process
serves only the first and then exits.

\subsection*{Exercise 8.2 --- Resource lifetime}

The listener exists from socket creation through the outer \code{with}. The OS
queues connection state after listen. Accept returns a new descriptor owned by
the inner \code{with}; closing it sends/causes TCP teardown independently of the
listener. Leaving either block closes its descriptor even during an exception.

\subsection*{Exercise 9.1 --- Fragment table}

For a request of length $n$, there are $n-1$ interior two-fragment divisions.
Before the second fragment, the parser returns incomplete because no full final
delimiter exists (unless the division is after the entire header). After both,
the accumulated bytes are identical, so the parsed result must be identical.
This property is ideal for a loop-driven unit test.

\subsection*{Exercise 9.2 --- Limit placement}

Request-line limit can be enforced when CRLF arrives or the buffer exceeds the
limit before CRLF. Header-byte limit applies while searching for CRLF--CRLF.
Field count applies after lines are available. Body limit applies immediately
after a valid declared length, before body bytes arrive. These early checks avoid
waiting or allocating for already-invalid messages.

\subsection*{Exercise 10.1 --- Adversarial table}

Valid examples should include GET, zero-length POST, repeated non-framing fields,
and percent-encoded Unicode. Incomplete examples include partial request line,
partial field delimiter, declared body short by one, and second request prefix.
Invalid examples include missing Host, duplicate length, transfer coding, and
invalid percent escape. Expected outcomes are Request, None, or HTTPError with
400/413/431/501 as appropriate.

\subsection*{Exercise 10.2 --- Invariant proof}

Before pop, the buffer is an unconsumed prefix of one or more stream messages.
After success, it contains exactly the unconsumed suffix after one complete
message. Header end identifies the body start; validated Content-Length identifies
the first byte after the body. Deleting their sum therefore removes neither too
few nor too many bytes.

\subsection*{Exercise 11.1 --- Serializer matrix}

GET/200: Content-Length equals body and body is sent. HEAD/200: same length but
zero transmitted content. GET/204 and GET/304: no content and this serializer
removes Content-Length. GET/500: problem JSON length and body are sent. All have
a status line and connection policy.

\subsection*{Exercise 11.2 --- Injection attempt}

\code{Headers({"X-Test": "ok\\r\\nInjected: yes"})} raises ValueError. Percent
encoding is representation-specific and recipients do not generally decode it
in field lines; applying it blindly changes legitimate values and still leaves
other control-character contexts. Validate the grammar at the field boundary.

\subsection*{Exercise 12.1 --- Route conflict table}

Static \code{/tasks/statistics} must beat \code{/tasks/\{task\_id\}}. GET on both
resolves; POST on the statistics path returns 405 if only GET exists; an unknown
two-segment path can match the parameter route; an unknown three-segment path is
404. For every matched path, allowed includes OPTIONS and includes HEAD when GET
exists.

\subsection*{Exercise 12.2 --- Nested resource}

Use \path{/projects/{project_id}/tasks} and
\path{/projects/{project_id}/tasks/{task_id}}. The router captures both
strings. Application code converts IDs, checks positivity, confirms project and
task existence, and ensures the task belongs to the project. Authorization also
belongs after routing.

\subsection*{Exercise 13.1 --- Error taxonomy}

HTTP line/framing errors belong to parser and produce 400/4xx/501 with close.
Malformed JSON belongs to Request JSON and is 400. Invalid fields belong to Task
validation and are 422. Missing task is resource logic and 404. SQL constraint
unexpectedly reached or programmer exception is 500 and logged. Only expected
public errors carry detailed client-safe messages.

\subsection*{Exercise 13.2 --- Request ID middleware}

Accept an incoming ID only from a trusted policy, for example 1--64 ASCII letters,
digits, hyphen, and underscore; otherwise generate a UUID. Put it in request
context and an \code{X-Request-ID} response field. Use structured logging so it
is a value, not format text. Do not accept CR/LF, unlimited length, or arbitrary
Unicode control characters.

\subsection*{Exercise 14.1 --- Capacity timeline}

Clients 1--2 occupy workers; 3--5 fill the queue; 6 receives 503/close. When one
worker times out, it closes its slow client and immediately takes client 3, leaving
two queued slots used. Exact TCP accept backlog state can exist before the
application queue and should be drawn separately.

\subsection*{Exercise 14.2 --- Concurrency test}

Start on port zero, wait for readiness, then use a 50-worker executor. Each client
sets a finite timeout, sends GET with Connection close, consumes the response,
and returns status. Collect with a global deadline. With the default 16/64
capacity, all 50 should be 200 locally; a deliberately smaller queue may allow
200 or 503, but never hangs or malformed responses.

\subsection*{Exercise 15.1 --- Add a route safely}

Register GET \code{/echo/\{value\}}. Read the captured value and
\code{request.query.get("tag", [])}; return both through Response.json. Tests
should include UTF-8 percent encoding, \code{?tag=} as an empty string, repeated
tags preserving order, HEAD with no transmitted body in integration, OPTIONS
with Allow, and POST returning 405.

\subsection*{Exercise 15.2 --- Boundary review}

Example parser contract: feed accepts bytes and may raise a limit error; pop
returns one Request, None, or HTTPError; its invariant is an unconsumed stream
prefix; it never discards following bytes. Compare with tests for fragmentation,
combined requests, framing conflicts, and limits. A useful missing case is close
during a declared body, tested at the server reader boundary.

\section{Backend and production bridge}

\subsection*{Exercise 16.1 --- Contract examples}

Examples must include Host, JSON media type, byte-correct length, status, response
media type, and body. Creation returns 201/Location; list returns the envelope;
individual reads/updates return task JSON; delete returns 204. Error fixtures
cover 400, 404, 405/Allow, 411, 413, 415, 422 field errors, and generic 500.

\subsection*{Exercise 16.2 --- completed\_at design}

Make it server-owned and nullable. Transition into done sets current UTC if null;
transition out clears it; repeated done assignment preserves or deliberately
refreshes it according to the written contract. Add a nullable schema column and
state/status consistency check if expressible. Include it in representations,
exclude client writes, migrate version, and preserve old clients that ignore it.

\subsection*{Exercise 17.1 --- Constraint failure}

Calling create with status \code{almost} raises an integrity error and the context
manager rolls back, leaving row count unchanged. HTTP validation remains necessary
because it produces a stable 422 with field detail; a raw database exception would
otherwise cross the generic 500 boundary and waste database work.

\subsection*{Exercise 17.2 --- Query plan}

An equality predicate on indexed status can use \code{tasks\_status\_idx}. A
pattern \code{\%word\%} cannot seek from an ordered B-tree prefix, so SQLite
normally scans rows. Full-text search or a different index strategy is needed at
scale; parameterization remains required either way.

\subsection*{Exercise 18.1 --- End-to-end extension}

Register GET \code{/tasks/summary} before or independently of parameter order;
static precedence ensures it wins. Repository SQL groups by status and fills
missing enum values with zero in application output. Test all-zero, mixed counts,
HEAD, OPTIONS, and wrong method. The handler returns one JSON object and performs
no identifier conversion.

\subsection*{Exercise 18.2 --- Optimistic concurrency design}

Generate a strong ETag from a version column or canonical representation. GET
returns it. PUT/PATCH requires or optionally accepts If-Match; a mismatch returns
412 Precondition Failed. The database update includes \code{WHERE id=? AND
version=?}, increments version atomically, and distinguishes missing resource
from stale version. Define policy for absent If-Match before implementation.

\subsection*{Exercise 19.1 --- Boundary expansion}

For a 1 MiB body limit, declare lengths $limit-1$, $limit$, and $limit+1$. The
first two remain acceptable if fully received; the last is rejected as soon as
fields are parsed, without receiving body bytes. Repeat at a small test limit to
avoid allocating large fixtures.

\subsection*{Exercise 19.2 --- Mini threat model}

Assets include task data, host capacity, database integrity, and logs. Entry
points are TCP bytes, HTTP fields/body/query/path, CLI/database path, and local
file access. Controls include loopback, strict parser, limits, parameterized SQL,
and safe 500. Major missing controls include identity/authorization, TLS, rate
limits, process isolation, backup, and operational monitoring. Public exposure
would make unauthorized data access and resource exhaustion top risks.

\subsection*{Exercise 20.1 --- Production comparison}

A sound comparison maps server protocol parsing to a production server; Request
and Response to framework types; route decorator and middleware to their mature
equivalents; worker/timeouts/shutdown to server settings; SQLite repository to
the framework's database layer. Capabilities absent here commonly include HTTP/2,
TLS, streaming, async, multipart forms, cookies, sessions, authentication,
dependency injection, schema validation, and deployment tooling.

\subsection*{Exercise 20.2 --- Capstone retrospective}

A strong note names specific prior beliefs, cites reproducible experiments, and
connects the most security-sensitive decision to framing agreement. It should
defend at least the boundary between production protocol server and application:
use maintained parsers/TLS/workers while keeping domain logic independently
tested. Personal examples vary; unsupported confidence without evidence does not.
